% ============================================================
% Section 70: 二维正方晶格的能隙与金属性条件
% ============================================================
\section{固体理论：二维正方晶格的能隙与金属性条件}
\label{sec:2d-square-gap}

\begin{modelbox}[题目：二维正方晶格的自由电子能量比、能隙与金属性条件]
考虑一个二维正方晶体。
\begin{enumerate}
    \item 位于其第一布里渊区顶角的自由电子能量是位于边界中点处自由电子能量的 $b$ 倍，求 $b$；
    \item 晶体势场为 $V(x,y)=-2V_0\bigl(\cos\frac{2\pi x}{a}+\cos\frac{2\pi y}{a}\bigr)$，近似求解第一布里渊区边界面中点的能隙；
    \item 设 (2) 中的结果是严格的，且材料是二价的，写出系统呈金属性的条件。
\end{enumerate}
\end{modelbox}

\begin{tipbox}[考点快速分析]
\begin{itemize}
    \item \textbf{自由电子能量比}：$E\propto k^2$，只需比较 $k$ 的模长
    \item \textbf{能隙来源}：近自由电子近似中，布里渊区边界处的能隙由势场对应倒格矢的傅里叶分量决定：$\Delta = 2|V_G|$
    \item \textbf{金属性判据}：二价材料恰好填满第一布里渊区；若能带交叠则为金属，分离则为绝缘体
    \item \textbf{能带交叠条件}：比较布里渊区顶角能量与相邻能带底部能量
\end{itemize}
\end{tipbox}

\begin{derivationbox}[标准解题过程]
\textbf{Step 1: 布里渊区顶角与边界中点的能量比}

二维正方晶格的倒格子也是正方晶格，倒格常数为 $2\pi/a$。第一布里渊区是边长为 $2\pi/a$ 的正方形（中心在原点）。

\begin{itemize}
    \item \textbf{边界中点}（$X$点）：如 $(\pi/a,0)$，到原点距离 $k_X = \pi/a$
    \item \textbf{顶角}（$M$点）：如 $(\pi/a,\pi/a)$，到原点距离 $k_M = \sqrt{(\pi/a)^2+(\pi/a)^2} = \sqrt{2}\,\pi/a$
\end{itemize}

自由电子能量 $E = \hbar^2k^2/2m \propto k^2$，故
\begin{equation}
b = \frac{E_M}{E_X} = \frac{k_M^2}{k_X^2} = \frac{2(\pi/a)^2}{(\pi/a)^2} = \boxed{2.}
\end{equation}

\textbf{Step 2: 第一布里渊区边界中点的能隙}

将势场展开为平面波的线性组合：
\begin{align}
V(x,y) &= -2V_0\Bigl(\cos\frac{2\pi x}{a}+\cos\frac{2\pi y}{a}\Bigr) \notag\\
&= -V_0\Bigl(e^{i\frac{2\pi}{a}x}+e^{-i\frac{2\pi}{a}x}+e^{i\frac{2\pi}{a}y}+e^{-i\frac{2\pi}{a}y}\Bigr).
\end{align}

第一布里渊区边界中点 $X$ 对应波矢 $\bm{k}=(\pi/a,0)$。在该点，自由电子平面波 $e^{i\bm{k}\cdot\bm{r}}$ 与 $e^{i(\bm{k}-\bm{G})\cdot\bm{r}}$ 简并，其中倒格矢 $\bm{G}=(2\pi/a,0)$ 恰好将 $k$ 映射到等价的 $-k$。

在近自由电子近似（简并微扰论）中，布里渊区边界上的能隙由势场中对应倒格矢的傅里叶分量决定：
\begin{equation}
\Delta_X = 2|V_G|.
\end{equation}
由势场展开式，对应 $\bm{G}=(2\pi/a,0)$ 的傅里叶系数为 $-V_0$，即 $|V_G|=|V_0|$。故
\begin{equation}
\boxed{\Delta_X = 2|V_0|.}
\end{equation}
（同理，$Y$ 点 $(0,\pi/a)$ 的能隙也是 $2|V_0|$。）

\textbf{Step 3: 二价材料呈金属性的条件}

二维正方晶格每个原胞含 1 个原子。二价材料每个原胞贡献 2 个价电子。考虑自旋简并，2 个电子恰好填满第一布里渊区（第一能带）。

\textbf{金属与绝缘体的判据}：
\begin{itemize}
    \item 若第一能带顶 $<$ 第二能带底，且费米能级落在能隙中 $\rightarrow$ \textbf{绝缘体}
    \item 若能带发生交叠（第一能带某些状态能量高于第二能带底部）$\rightarrow$ \textbf{金属}
\end{itemize}

近自由电子近似下：
\begin{itemize}
    \item 第二能带底部在 $X$ 点：$E_{\text{底}}^{(2)} \approx E_X + |V_0| = \dfrac{\hbar^2}{2m}\Bigl(\dfrac{\pi}{a}\Bigr)^2 + |V_0|$
    \item 第一能带中能量较高的自由电子态在 $M$ 点：$E_M = \dfrac{\hbar^2}{2m}\cdot 2\Bigl(\dfrac{\pi}{a}\Bigr)^2 = 2E_X$
\end{itemize}

在简化模型中（将布里渊区边界处的能隙整体效果纳入比较），能带交叠的判据为：
\begin{equation}
E_X + \Delta_X \le E_M,
\end{equation}
即第一能带在 $X$ 方向的顶部（自由电子能量 $E_X$）加上该方向完整能隙 $\Delta_X=2|V_0|$ 后，不超过 $M$ 点能量。代入得
\begin{equation}
\frac{\hbar^2}{2m}\Bigl(\frac{\pi}{a}\Bigr)^2 + 2|V_0| \le \frac{\hbar^2}{m}\Bigl(\frac{\pi}{a}\Bigr)^2.
\end{equation}
化简：
\begin{equation}
2|V_0| \le \frac{\hbar^2}{2m}\Bigl(\frac{\pi}{a}\Bigr)^2
\quad\Longrightarrow\quad
\boxed{|V_0| \le \frac{\pi^2\hbar^2}{4ma^2}.}
\end{equation}

\textbf{结论}：当周期势振幅 $|V_0|$ 小于临界值 $\pi^2\hbar^2/(4ma^2)$ 时，能带发生交叠，二价材料呈现\textbf{金属性}；反之，若势场足够强（能隙足够大），能带分离，材料为\textbf{绝缘体}。
\end{derivationbox}

\begin{formulabox}[答案汇总]
\begin{center}
\begin{tabular}{ll}
\toprule
\textbf{问题} & \textbf{答案} \\
\midrule
(1) 能量比值 $b$ & $b=2$ \\
(2) 边界中点能隙 & $\Delta_X = 2|V_0|$ \\
(3) 金属性条件 & $|V_0| \le \pi^2\hbar^2/(4ma^2)$ \\
\bottomrule
\end{tabular}
\end{center}
\end{formulabox}

\begin{insightbox}[物理图像]
\begin{itemize}
    \item 二维正方晶格第一布里渊区顶角 $M$ 点到原点的距离是边界中点 $X$ 点的 $\sqrt{2}$ 倍，故自由电子能量比值为 $2$。
    \item 近自由电子近似下，布里渊区边界上的能隙由势场对应倒格矢的傅里叶分量决定，这是能带论中最核心的结论之一。
    \item 二价材料是金属还是绝缘体，本质上取决于第一、二能带是否交叠。\textbf{弱周期势}（小 $V_0$）有利于能带交叠，形成金属——这正是碱土金属（如 Mg、Ca）虽然是二价但仍为金属的原因。
\end{itemize}
\end{insightbox}
